\chapter{Theory}
\label{ch:theory}

This chapter summarises what \pycmf{} computes, in the notation of the code.
The derivations are given in the papers listed in
Appendix~\ref{ch:references}.

\section{OBMP2}
\label{sec:theory-obmp2}

\subsection{Correlated Fock operator}

The electronic Hamiltonian $\hat H$ is transformed with an anti-Hermitian
doubles operator $\hat A$ built from first-order amplitudes. The
Baker--Campbell--Hausdorff (BCH) expansion is truncated at second order, with
the Fock operator $\hat F$ in the double commutator, and every commutator is
reduced to one-body operators and constants by a cumulant approximation,
denoted $[\ldots]_1$:
\begin{equation}
  \hat H_{\mathrm{OBMP2}} = \hat H_{\mathrm{HF}} + \big[\hat H,\hat A\big]_1
  + \tfrac12\big[[\hat F,\hat A],\hat A\big]_1
  = \hat{\bar F} + \bar C ,
  \qquad \bar f_{pq} = f_{pq} + v_{pq}.
  \label{eq:hobmp2}
\end{equation}
Here $f$ is the Hartree--Fock (HF) Fock matrix, $v$ is the
\emph{correlation potential} and $\fbar$ the \emph{correlated Fock matrix}.
The orbitals are the eigenvectors of $\fbar$ and the orbital energies
$\varepsilon_p$ its eigenvalues.

\subsection{Amplitudes and self-consistency}

The amplitudes are rebuilt from the current orbitals and orbital energies,
\begin{equation}
  t^{ab}_{ij} = c\,
  \frac{(ia|jb)}{\varepsilon_i+\varepsilon_j-\varepsilon_a-\varepsilon_b-\Delta},
  \label{eq:amplitudes}
\end{equation}
where $c$ is a scaling factor ($c=1$ in OBMP2) and $\Delta$ an optional level
shift (attribute \code{level_shift}). The cycle
\[
  \text{orbitals} \;\to\; t \;\to\; v \;\to\; \fbar \;\to\; \text{new orbitals}
\]
is repeated until the energy is stable and the Brillouin-type condition
$\bar f_{ia}=0$ holds. Starting from HF orbitals, the energy of the first
cycle is the conventional MP2 energy.

\subsection{Energy}

The energy of a cycle is
\begin{equation}
  E = E_{\mathrm{HF}}[\Phi] + c_0 + \sum_{i}^{\mathrm{occ}} v_{ii},
  \label{eq:energy}
\end{equation}
where $E_{\mathrm{HF}}[\Phi]$ is the HF energy of the current determinant and
$c_0$ the constant of the correlation potential. Because the amplitudes
depend on the orbital energies, this functional is not stationary with respect
to orbital rotations: the converged orbitals are defined by the eigenvalue
problem, not by a variational principle. The error of the energy is therefore
of first order in the residual $\bar f_{ia}$, and a tight convergence
threshold on $\bar f_{ia}$ (attribute \code{conv_tol_grad}) is needed when
energy differences below $10^{-6}$~Hartree matter.

\subsection{Spin blocks}

Amplitudes are stored in chemists' order, \code{t[i,a,j,b]}, as a list of
spin blocks $(x,y)$: the pair $(i,a)$ has spin $x$ and the pair $(j,b)$ has
spin $y$. The antisymmetrised amplitude is $\bar t = t - t^{\mathrm{exch}}$
for same-spin blocks and $\bar t = t$ for opposite-spin blocks. An
unrestricted reference has the four blocks $\alpha\alpha$, $\beta\beta$,
$\alpha\beta$ and $\beta\alpha$. A restricted reference has one spin channel
and two blocks (same spin, opposite spin); the spin-flipped partners
contribute a factor of two to the constants. The same routines
(\code{pycmf.mol.bch}) therefore serve both cases.

\subsection{Working equations}

With $f$ the HF Fock matrix in the current orbital basis and summation over
repeated indices, each block $(x,y)$ adds to the unsymmetrised matrices
$c_1^{x}$, $c_1^{y}$ and to the constant $c_0$.

\paragraph{First order, $[\hat H,\hat A]_1$.}
\begin{align}
  c_1^{y}[p,k] &\mathrel{+}= 2\,(ia|pl)\,\bar t_{ia,kl}, &
  c_1^{y}[p,l] &\mathrel{-}= 2\,(ia|kp)\,\bar t_{ia,kl}, \\
  c_1^{y}[j,b] &\mathrel{+}= 2\,f^{x}_{ia}\,\bar t_{ia,jb}, &
  c_0 &\mathrel{-}= (ia|jb)\,\bar t_{ia,jb}.
\end{align}

\paragraph{Second order, $\tfrac12[[\hat F,\hat A],\hat A]_1$.}
With the intermediates $V_{ia,jd} = f^{x}_{ac}\,\bar t_{ic,jd}$ and
$O_{ia,jd} = f^{x}_{ik}\,\bar t_{ka,jd}$,
\begin{align}
  c_1^{y}[j,k] &\mathrel{+}= t_{ia,jb}\,(V-O)_{ia,kb}, &
  c_1^{x}[i,k] &\mathrel{+}= t_{ia,jb}\,V_{ka,jb}, \\
  c_1^{x}[a,c] &\mathrel{+}= t_{ia,jb}\,O_{ic,jb}, &
  c_1^{y}[b,d] &\mathrel{+}= t_{ia,jb}\,(O-V)_{ia,jd}, \\
  c_0 &\mathrel{+}= t_{ia,jb}\,(O-V)_{ia,jb}, &
  c_1^{x}[i,a] &\mathrel{+}= \bar t_{ia,jb}\,y^{y}_{jb},
\end{align}
with $y^{y}_{jb}=\sum_{x'} f^{x'}_{kc}\,\bar t^{\,x'y}_{kc,jb}$, and, with
$R^{\mathrm{occ}}_{ik} = t_{ia,jb}\bar t_{ka,jb}$ and
$R^{\mathrm{vir}}_{ac} = t_{ia,jb}\bar t_{ic,jb}$ accumulated over the blocks
with first spin $x$,
\begin{equation}
  c_1^{x}[p,i] \mathrel{-}= f^{x}_{pk}R^{\mathrm{occ}}_{ik},\qquad
  c_1^{x}[p,c] \mathrel{-}= f^{x}_{pa}R^{\mathrm{vir}}_{ac}.
\end{equation}
Finally $v = \tfrac12\,(c_1 + c_1^{\mathsf T})$. The second-order part needs
no two-electron integrals.

\paragraph{First-order variant.}
\code{bch_order=1} keeps only the first commutator with half weight,
$\hat H + \tfrac12[\hat H,\hat A]_1$. It was the default of some modules of
version~0.2; the second-order form (\code{bch_order=2}) is the default now.

\subsection{Cost}

One cycle needs the integrals $(ia|jb)$, $(ia|pb)$ and $(ia|jp)$ and
contractions that scale as $o^2v^2N$ (first order) and $o^2v^3$ (second
order), where $o$, $v$ and $N$ are the numbers of occupied, virtual and all
orbitals. The cost is therefore $O(N^5)$ per cycle, like MP2. With exact
integrals the first index pair is transformed once per cycle,
$(ia|\mu\nu)$; with density fitting the three-index tensors $B^Q_{pa}$ and
$B^Q_{ip}$ are used. Memory is dominated by the amplitudes, $O(o^2v^2)$, and
by the three-index or half-transformed integrals.

\section{O2BMP2}
\label{sec:theory-o2bmp2}

Same-spin amplitudes are dropped and opposite-spin amplitudes are scaled,
$c_{\mathrm{SS}}=0$ and $c_{\mathrm{OS}}=c$. The scaling acts on the amplitudes
that enter $v$, so it changes the orbitals and orbital energies and not only
the energy. More generally, the attributes \code{c_os} and \code{c_ss} scale
the opposite-spin and same-spin amplitudes independently; $(1,1)$ is OBMP2,
$(0,0)$ is Hartree--Fock.

There is no universal value of $c_{\mathrm{OS}}$. The published choices are
1.2 (singlet--triplet gaps, charge-transfer excitations), 1.1 (ionization
potentials) and 1.7 (inverted singlet--triplet gaps). \code{pycmf.O2BMP2}
therefore has no default. An $O(N^4)$ algorithm based on the Laplace
transformation is not implemented.

\section{One-body double hybrids: OBDH and O2BDH}
\label{sec:theory-obdh}

For a double-hybrid functional with exact-exchange fraction $a_x$ and
correlation weight $a_c$,
\begin{equation}
  \fbar = f^{\mathrm{KS}}\big[a_x\,\mathrm{HF} + (1-a_x)\,X,\ (1-a_c)\,C\big]
          + a_c\,v_{\mathrm{OBMP2}},
  \qquad
  E = E^{\mathrm{KS}} + a_c\big(c_0 + \textstyle\sum_i v_{ii}\big),
\end{equation}
where $X$ and $C$ are the semilocal exchange and correlation functionals.
$v_{\mathrm{OBMP2}}$ is built with the HF Fock matrix of the current orbitals
and with amplitudes whose denominators are the eigenvalues of $\fbar$. With
$a_x=a_c=1$ the method is OBMP2; with $a_c=0$ it is the hybrid functional.
O2BDH uses opposite-spin amplitudes only, scaled by $c_{\mathrm{OS}}$.

The Kohn--Sham part is evaluated by \pyscf{} (libxc). The default functional
is B2PLYP. The global double hybrids of Table~\ref{tab:dh} are available by
name. Their parameters were fitted with a perturbative second-order energy on
Kohn--Sham orbitals and have not been re-optimised for the self-consistent
scheme.

\begin{table}[!htb]
\centering
\small
\caption{Double hybrids available by name (\code{pycmf.mol.obdh.FUNCTIONALS}).}
\label{tab:dh}
\begin{tabular}{@{}l l l l l l@{}}
\toprule
Name & $a_x$ & $a_c$ & Exchange & Correlation & Reference \\
\midrule
B2PLYP    & 0.53 & 0.27  & B88   & LYP & Grimme (2006) \\
mPW2PLYP  & 0.55 & 0.25  & mPW91 & LYP & Schwabe and Grimme (2006) \\
B2T-PLYP  & 0.60 & 0.31  & B88   & LYP & Tarnopolsky \emph{et al.} (2008) \\
B2GP-PLYP & 0.65 & 0.36  & B88   & LYP & Karton \emph{et al.} (2008) \\
B2K-PLYP  & 0.72 & 0.42  & B88   & LYP & Tarnopolsky \emph{et al.} (2008) \\
PBE0-DH   & 0.5  & 0.125 & PBE   & PBE & Br\'emond and Adamo (2011) \\
PBE-QIDH  & $3^{-1/3}$ & $1/3$ & PBE & PBE & Br\'emond \emph{et al.} (2014) \\
PBE0-2    & $2^{-1/3}$ & $1/2$ & PBE & PBE & Chai and Mao (2012) \\
\bottomrule
\end{tabular}
\end{table}

OBDH and O2BDH have not yet been described in a publication; this manual and
the docstring of \code{pycmf.mol.obdh} are their definition.

\section{Orbital-optimized double hybrids (OO-DH)}
\label{sec:theory-oodh}

OO-DH minimizes the double-hybrid energy itself with respect to the orbitals
(Peverati and Head-Gordon; OO-RI-B2PLYP of Neese \emph{et al.}):
\begin{equation}
  E[C] = E^{\mathrm{KS}}\big[a_x\,\mathrm{HF} + (1-a_x)\,X,\ (1-a_c)\,C;\ D\big]
         + a_c\,E_2[C],
\end{equation}
where $E_2$ is the MP2 energy built with the Kohn--Sham Fock matrix $f$ of the
same orbitals in semicanonical form (the occupied--occupied and
virtual--virtual blocks of $f$ are diagonalized, the singles are left out).
This is the OO-MP2 Lagrangian of Bozkaya \emph{et al.} with a Kohn--Sham
reference. In spin orbitals the Hylleraas functional is
\begin{equation}
  E_2 = \tfrac12\sum_{ijab} t_{ij}^{ab}\langle ij\|ab\rangle
        + \sum_{pq} f_{pq}\gamma_{pq},\qquad
  \gamma_{ij} = -\tfrac12\sum_{kab}t_{ik}^{ab}t_{jk}^{ab},\quad
  \gamma_{ab} = \tfrac12\sum_{ijc}t_{ij}^{ac}t_{ij}^{bc},
\end{equation}
and it is stationary in $t$. For the rotation
$\phi_i\to\phi_i+x_{ai}\phi_a$ the exact orbital gradient is therefore
\begin{equation}
  \frac{\partial E}{\partial x_{ai}} = 2f_{ai} + a_c\Big[
    \sum_{jbc}t_{ij}^{bc}\langle aj\|bc\rangle
  - \sum_{jkb}t_{jk}^{ab}\langle jk\|ib\rangle
  + 2\,(f\gamma-\gamma f)_{ai} + 2\,R[\gamma]_{ai}\Big],
\end{equation}
with $R[\gamma]$ the response of the Kohn--Sham Fock matrix to the density
$\gamma$ (Coulomb, scaled exact exchange and the exchange--correlation
kernel). The last term is what makes the gradient exact for a functional of
the density. The self-consistent driver diagonalizes
\begin{equation}
  \fbar = f + a_c v,\qquad
  v_{ai} = \frac{1}{2a_c}\Big(\frac{\partial E}{\partial x_{ai}} - 2f_{ai}\Big),
  \qquad v_{ij}=v_{ab}=0,
\end{equation}
so that the converged $\bar f_{ai}$ is half the orbital gradient. Because the
energy is stationary in the orbitals, $D + a_c\gamma$ is the relaxed density:
its expectation value of a one-electron operator is the derivative of the
energy (no $Z$-vector equations).

OO-DH differs from OBDH (Section~\ref{sec:theory-obdh}): OBDH adds the OBMP2
correlation potential to the Kohn--Sham Fock matrix and its energy is not
variational; OO-DH uses the exact gradient of the double-hybrid energy.
With $a_x=a_c=1$ and no semilocal parts it is OO-MP2; with $a_c=0$ it is the
hybrid functional. $c_{\mathrm{OS}}$ and $c_{\mathrm{SS}}$ scale the two spin
components of $E_2$.

\section{State-specific excited states}
\label{sec:theory-mom}

The correlated Fock operator is a one-body operator, so non-Aufbau
determinants can be optimised as in $\Delta$SCF. The calculation starts from a
MOM-HF determinant. After each diagonalisation the orbitals with the largest
overlap with the occupied space of the previous cycle (\code{'mom'}, Gilbert,
Besley and Gill) or of the initial determinant (\code{'imom'}, Barca, Gilbert
and Gill) are occupied. Excitation energies are differences of total
energies. An open-shell singlet obtained this way is a spin-mixed determinant
($\langle S^2\rangle\approx 1$).

\section{Transition moments and NOCI}
\label{sec:theory-noci}

The ground and excited determinants of a state-specific calculation have
different, non-orthogonal orbitals. For each spin the generalized
Slater--Condon rules are used in the corresponding-orbital form,
\begin{equation}
  X = C_L^{\mathsf T} S\, C_R = U\,\mathrm{diag}(\sigma)\,V^{\mathsf T},\qquad
  s = \det X,\qquad
  P = \pm\sum_i \Big[\prod_{k\neq i}\sigma_k\Big]\,|r_i\rangle\langle l_i|,
\end{equation}
with $l=C_LU$ and $r=C_RV$. The cofactor weights avoid a division by the
overlap, so the expressions also hold when $X$ is singular. The two-electron
terms use products of the same cofactors (Lemke, Kussmann and Ochsenfeld).

\paragraph{Transition dipole.} The two determinants are first made orthogonal
by the symmetric (L\"owdin) transformation (Bourne Worster, Feighan and
Manby), then $\boldsymbol\mu$ is taken between them and
$f = \tfrac23\,\Delta E\,|\boldsymbol\mu|^2$. This is a determinant-level
moment: the correlation enters through the orbitals only. For a closed-shell
molecule the excited determinant is spin mixed; \code{singlet=True}
multiplies $\boldsymbol\mu$ by $\sqrt2$, assuming the $\alpha$ and $\beta$
branches are equivalent and orthogonal.

\paragraph{NOCI.} The ground-state determinant and state-specific excited
determinants (one per excitation channel and spin) are combined by solving
$HC = SCE$ after canonical orthogonalization of $S$. Three Hamiltonians are
available:
\begin{itemize}
  \item \code{'bare'}: the physical Hamiltonian; the diagonal is the HF energy
        of each OBMP2 determinant.
  \item \code{'overlap_symmetric'}:
        $H_{IJ} = H^{\mathrm{bare}}_{IJ} + \tfrac12(\Delta_I+\Delta_J)S_{IJ}$
        with $\Delta_I = E_I - H^{\mathrm{bare}}_{II}$, so the diagonal is the
        OBMP2 energy.
  \item \code{'state_specific_fock'} (default): each determinant carries its
        correlated Fock operator $F_I$ and a constant $c_I$ with
        $c_I+\sum_{\mathrm{occ}}F_I=E_I$;
        $H_{IJ}=\tfrac12(c_I+c_J)S_{IJ}+\sum_\sigma\mathrm{tr}
        [\tfrac12(F_{I\sigma}+F_{J\sigma})D_{IJ\sigma}]$.
\end{itemize}
The last two are model choices; only their diagonal is fixed by the theory.
The averaged Fock operator is a one-body operator, so it does not couple the
$\alpha$ and $\beta$ determinants of one channel, which differ by two spin
orbitals (known issue~8).

\section{One-body downfolding (OBDF)}
\label{sec:theory-obdf}

Let $T_{\mathrm{act}}$ be the amplitudes with all four indices in an active
space. The external correlation potential and constant are
\begin{equation}
  v^{\mathrm{ext}} = v[T] - v[T_{\mathrm{act}}],\qquad
  c_0^{\mathrm{ext}} = c_0[T]-c_0[T_{\mathrm{act}}],
\end{equation}
and the active-space Hamiltonian is
\begin{equation}
  \hat H_{\mathrm{act}} = E_{\mathrm{core}}
  + \sum_{tu}\bar h_{tu}\,\hat a^\dagger_t\hat a_u
  + \tfrac12\sum_{tuvw}(tu|vw)\,\hat a^\dagger_t\hat a^\dagger_v\hat a_w\hat a_u,
\end{equation}
\begin{equation}
  \bar h = h + v^{\mathrm{HF}}_{\mathrm{core}} + v^{\mathrm{ext}},\qquad
  E_{\mathrm{core}} = E_{\mathrm{nuc}} + E^{\mathrm{HF}}_{\mathrm{core}}
  + c_0^{\mathrm{ext}} + \sum_{i\in\mathrm{core}} v^{\mathrm{ext}}_{ii}.
\end{equation}
Only the one-body integrals are dressed; the two-electron integrals are the
bare ones, so $\hat H_{\mathrm{act}}$ has the structure of a CASCI
Hamiltonian and can be handed to any solver (FCI, DMRG, selected CI, quantum
algorithms).

Two limits hold exactly and are tested: if the active space contains all
orbitals, $\hat H_{\mathrm{act}}=\hat H$ (FCI energy); if it is empty,
$E_{\mathrm{core}}=E_{\mathrm{OBMP2}}$. For any active space the reference
determinant has the energy
\begin{equation}
  \langle\Phi|\hat H_{\mathrm{act}}|\Phi\rangle
  = E_{\mathrm{OBMP2}}-E_c[T_{\mathrm{act}}],\qquad
  E_c[T_{\mathrm{act}}]=c_0[T_{\mathrm{act}}]+\sum_i v[T_{\mathrm{act}}]_{ii}:
\end{equation}
the correlation energy of the active double excitations is what is left to
the active-space solver.

\section{Properties}
\label{sec:theory-properties}

\begin{itemize}
  \item \emph{HF-like density}: the density of the OBMP2 determinant.
  \item \emph{MP2-like density}: $\gamma_{\mathrm{oo}}=1-R^{\mathrm{occ}}$ and
        $\gamma_{\mathrm{vv}}=R^{\mathrm{vir}}$ with the converged amplitudes
        (unrelaxed). After one cycle from HF orbitals it equals the unrelaxed
        MP2 density.
  \item $\langle S^2\rangle$ of the determinant, optionally with the
        first-order doubles correction
        $\langle S^2\rangle = \langle S^2\rangle_{\mathrm{det}}
        - \sum_{ia,jb} S_{ib}S_{ja}t_{ia,jb}$, where $S$ is the overlap of
        the $\alpha$ and $\beta$ orbitals.
  \item Dipole moments and isotropic hyperfine coupling constants as
        expectation values over either density.
  \item Koopmans-type ionization potentials and electron affinities from the
        eigenvalues of $\fbar$.
\end{itemize}

\section{Projection embedding in DFT}
\label{sec:theory-embed}

In the projection-based embedding (\code{pycmf.embed}, experimental) the two
levels have their own functionals. A Kohn--Sham calculation of the full
system with the low-level functional \code{xc_env} gives the densities
$\gamma_A$ and $\gamma_B$ of the subsystem and of the environment (SPADE
partition) and the embedding potential
\begin{equation}
  v_{\mathrm{emb}} = v^{\mathrm{env}}[\gamma_A+\gamma_B]-v^{\mathrm{env}}[\gamma_A]
  +\mu\,S\gamma_BS .
\end{equation}
The subsystem is then treated with OBMP2 or OBDH (with its own double hybrid)
in this potential, optionally in a virtual space truncated by concentric
localization, and the total energy is
\begin{equation}
  E = E^{\mathrm{high}}[\gamma_A'] + E^{\mathrm{env}}[\gamma_A'+\gamma_B]
  -E^{\mathrm{env}}[\gamma_A'] + \mu\,\mathrm{tr}(\gamma_A'S\gamma_BS),
\end{equation}
with $\gamma_A'$ the relaxed density of the embedded subsystem. Two limits are
exact and tested: with all atoms in the subsystem $E$ is the energy of the
high-level method, and with the same functional on both levels and $a_c=0$ it
is the Kohn--Sham energy of the full system.

\section{Periodic systems}
\label{sec:theory-pbc}

For a crystal with a $k$-point mesh of $\Nk$ points the working equations are
those of Section~\ref{sec:theory-obmp2}, spin-summed and with conservation
of crystal momentum. For the four $k$-points $(k_i,k_a,k_j,k_b)$ of a double
excitation $ia\to jb$,
\begin{equation}
  g_{iajb} = \frac{(ia|jb)}{\Nk},\qquad
  t_{iajb} = \frac{g^*_{iajb}}{\varepsilon_i-\varepsilon_a+\varepsilon_j-\varepsilon_b-\Delta},\qquad
  \bar t_{iajb} = \alpha\,t_{iajb} - \beta\,t_{ibja},
\end{equation}
with $\alpha = (c_{\mathrm{SS}}+c_{\mathrm{OS}})/2$ and $\beta =
c_{\mathrm{SS}}/2$ in the terms linear in the amplitudes, and the squares of
the scaling factors in the quadratic terms. All two-electron integrals are
products of Gaussian density-fitting tensors,
\begin{equation}
  (p k_1, q k_2 \,|\, r k_3, s k_4) = \sum_Q L^Q_{pq}(k_1,k_2)\,L^Q_{rs}(k_3,k_4),
\end{equation}
so no four-index transformation is performed. The algorithm, the use of
$k$-point symmetry and the quasiparticle methods built on KOBMP2 are
described in Chapter~\ref{ch:periodic}.
